% =====================================================================
%  Pakete und Layout – für pdflatex. Reihenfolge ist teilweise wichtig:
%  hyperref (fast) zuletzt, cleveref/bookmark danach.
% =====================================================================

% ---------- Zeichensatz & Sprache ----------
\usepackage[T1]{fontenc}
\usepackage[utf8]{inputenc}
\usepackage[english,\arbeitSprache]{babel}   % letzte Sprache = Hauptsprache
\usepackage{lmodern}
\usepackage{microtype}
\usepackage[autostyle=true,german=quotes]{csquotes}  % \enquote{…} → „…“

% ---------- Seitenlayout ----------
\usepackage[left=3cm,right=2.5cm,top=2.5cm,bottom=2.5cm,headsep=0.8cm,marginparwidth=2.1cm,marginparsep=0.25cm]{geometry}
\usepackage{setspace}
\onehalfspacing
\usepackage[automark,headsepline]{scrlayer-scrpage}
\clearpairofpagestyles
\ihead{\leftmark}
\ohead{\pagemark}
\pagestyle{scrheadings}
\setcounter{secnumdepth}{3}   % bis \subsubsection nummerieren
\setcounter{tocdepth}{2}      % Inhaltsverzeichnis bis \subsection

% ---------- Mathematik ----------
\usepackage{amsmath,amssymb}

% ---------- Abbildungen & Tabellen ----------
\usepackage{graphicx}
\graphicspath{{abbildungen/}}
\usepackage{float}
\usepackage[font=small,labelfont=bf,margin=0.5cm]{caption}
\usepackage{subcaption}
\usepackage{booktabs}
\usepackage{tabularx}
\usepackage{longtable}
\usepackage{multirow}
\usepackage{array}
\usepackage[table]{xcolor}
\usepackage{placeins}          % \FloatBarrier: Gleitobjekte vor dem nächsten Abschnitt ausgeben
\usepackage{pdfpages}          % \includepdf für Anhänge (Fragebögen, Datenblätter)

% ---------- Listen & Textbausteine ----------
\usepackage{enumitem}
\setlist{nosep}
\usepackage{url}
\usepackage{xspace}

% ---------- Quellcode ----------
\usepackage{listings}
\definecolor{codebg}{RGB}{247,249,248}
\definecolor{codekw}{RGB}{15,111,43}
\definecolor{codecomment}{RGB}{120,130,125}
\definecolor{codestring}{RGB}{11,110,153}
\lstset{
  basicstyle=\ttfamily\small,
  keywordstyle=\color{codekw}\bfseries,
  commentstyle=\color{codecomment}\itshape,
  stringstyle=\color{codestring},
  backgroundcolor=\color{codebg},
  numbers=left, numberstyle=\tiny\color{codecomment}, numbersep=8pt,
  frame=single, framerule=0.3pt, rulecolor=\color{gray!40},
  breaklines=true, breakatwhitespace=true,
  showstringspaces=false, tabsize=2, captionpos=b,
  inputencoding=utf8, extendedchars=true,
  literate={ä}{{\"a}}1 {ö}{{\"o}}1 {ü}{{\"u}}1 {Ä}{{\"A}}1 {Ö}{{\"O}}1 {Ü}{{\"U}}1 {ß}{{\ss}}1 {–}{{--}}1 {„}{{,,}}1 {“}{{''}}1,
}
\renewcommand{\lstlistingname}{Listing}
\renewcommand{\lstlistlistingname}{Listingverzeichnis}

% ---------- Abkürzungen ----------
% \ac{KEY} beim ersten Vorkommen ausgeschrieben, danach nur die Abkürzung.
% Liste in abkuerzungen.tex pflegen.
\usepackage[printonlyused]{acronym}

% ---------- Notizen während des Schreibens ----------
% \todo{…} am Rand, \todo[inline]{…} im Text, \missingfigure{…} als Platzhalter.
% Vor der Abgabe: Option "disable" setzen → \usepackage[disable]{todonotes}
\usepackage[colorinlistoftodos,prependcaption,textsize=footnotesize]{todonotes}

% ---------- Literatur ----------
\usepackage[
  backend=biber,
  style=\arbeitZitierstil,
  maxbibnames=99,
  maxcitenames=2,
  doi=true, isbn=false, url=true,
  date=year,
]{biblatex}
\DefineBibliographyStrings{ngerman}{andothers={et\,al\adddot}}

% ---------- Verweise (immer zuletzt) ----------
\usepackage[
  colorlinks=true,
  linkcolor=black, citecolor=black, urlcolor=blue!50!black,
  pdftitle={\arbeitTitel},
  pdfauthor={\arbeitAutor},
  pdfsubject={\arbeitArt},
  bookmarksnumbered=true,
  pdfstartview=FitH,
]{hyperref}
\usepackage{bookmark}

% ---------- Eigene Befehle ----------
\newcommand{\zB}{z.\,B.\xspace}
\newcommand{\dasheisst}{d.\,h.\xspace}   % (\dh ist in T1 bereits der Buchstabe ð)
\newcommand{\bzw}{bzw.\xspace}
\newcommand{\vgl}{vgl.\xspace}
\newcommand{\code}[1]{\texttt{#1}}
